\documentclass[11pt]{article}
\usepackage{icmlko}
\begin{document}

\title{늙은 축은 부호로 재면 안 보인다 --- 애니가 검색을 버리는 이유}
\author{월드모델 실험실}
\date{노트 492 · 2026년 8월 3일}
\maketitle

\begin{abstract}
노트 490 이 검색 시계열에 판 $-0.0283$ 을 매기면서 예외를 하나 남겼다 ---
\textbf{애니 $+0.0328$}. 열한 도메인 중 유일하게 큰 양수다(도서
$-0.1519$ · 아이돌 $-0.1170$ · 게임 $-0.0833$). \emph{애니는 검색
시계열이 없는 편이 낫다.}

설명 후보 셋을 차례로 쟀다.

\begin{center}
\begin{tabular}{lll}
\hline
후보 & 결과 & 근거 \\
\hline
축의 부호가 반대다 & \textbf{반증} & 애니 $+0.1286$ · 전체 $+0.1260$ \\
계절을 다시 적는다 & \textbf{반증} & 계절과 $+0.0408$($6$ 중 $4$ 등) \\
축이 늙었다 & \textbf{통과} & 아래 \\
\hline
\end{tabular}
\end{center}

그런데 셋째를 \emph{처음엔 틀리게} 쟀다. 부호 있는 상관으로 학습 대 유보를
견주면 아이돌이 $-0.0775$ 로 애니($-0.0681$)보다 더 늙어 보이는데,
아이돌에서는 시계열을 빼면 $-0.1170$ 을 잃는다. 설명이 안 선다.

\textbf{세기($|rho|$)로 재면 갈린다.} 학습에서 유보로 넘어가며 트렌드 축이
약해지는 도메인은 \textbf{애니 하나}다($-0.0592$). 나머지는 전부
\emph{강해진다} --- 도서 $+0.1812$ · 세계애니 $+0.1353$ · 게임 $+0.0833$
· 모바일 $+0.0530$ · 아이돌 $+0.0100$(웹툰만 $-0.0204$ 로 애니의 $1/3$).
\end{abstract}

\begin{figure}[t]
\centering
\includegraphics[width=\textwidth]{figs/f1.pdf}
\caption{\textbf{가} --- 학습 $\to$ 유보. \textbf{나} --- 대조를 나란히.
\textbf{다} --- 부호로 재면 둘이 같아 보인다.}
\end{figure}

\section{대조가 없었으면 못 말한다}

``애니 유보에서 트렌드가 약해진다''만으로는 부족하다 --- 애니 유보가 그냥
어려운 시기일 수도 있다. 그래서 같은 행에서 \textbf{기본 축 다섯}을 같이
쟀다(노트 484 조항 ②).

\begin{center}
\begin{tabular}{lrr}
\hline
도메인 & 트렌드 축 & 기본 축 다섯 \\
\hline
\textbf{애니} & $\mathbf{-0.0592}$ & $\mathbf{-0.0234}$ \\
웹툰 & $-0.0204$ & $-0.0361$ \\
아이돌 & $+0.0100$ & $-0.0760$ \\
모바일 & $+0.0530$ & $-0.0777$ \\
게임 & $+0.0833$ & $+0.1962$ \\
세계애니 & $+0.1353$ & $-0.0845$ \\
도서 & $+0.1812$ & $-0.0401$ \\
\hline
\end{tabular}
\end{center}

\textbf{애니의 기본 축은 거의 안 변한다}($-0.0234$ 로 일곱 중 둘째로
작다). 유보가 어려운 게 아니라 \textbf{트렌드가 특정해서 늙었다.} 다른
도메인들은 정반대 --- 기본 축이 약해지는데 트렌드는 강해진다(아이돌 ·
모바일 · 세계애니 · 도서 넷 다).

그러니 모형이 하는 일은 이렇다. 애니 학습 $1{,}567$ 행에서 $|rho| =
0.1340$ 인 축을 배워서, 그 축이 $0.0747$ 밖에 안 되는 유보 $585$ 행에
쓴다. 같은 축이 다른 도메인에서는 유보에서 \emph{더} 세지므로 공유 갈래는
그 방향으로 굳는다. \textbf{애니만 그 굳음을 뒤집어쓴다.}

\section*{조항}

\begin{quote}
\textbf{축이 늙었는지는 $|rho|$ 로 잰다.} 부호 있는 상관으로 재면 방향이
\emph{일관되게} 반대인 축(아이돌의 트렌드)을 늙은 것으로 잘못 읽는다 ---
그런 축은 모형이 배울 수 있으므로 안 늙은 것이다. \textbf{늙음은 방향이
아니라 세기의 문제다.}
\end{quote}

\section*{남은 것}

\begin{itemize}
\item \textbf{왜 애니에서만 늙었나} --- 이 노트는 \emph{늙었다}까지고
\emph{왜}는 아니다. 애니 검색은 $2025$ 년에 무엇이 달라졌나(노트 483
조항: 못 대면 이름 붙이기다).
\item \textbf{애니만 트렌드를 빼면 판이 오르나} --- 도메인 하나만 빼는
것은 지금 구조로는 못 한다(공유 열이다). 노트 449 의 ``한 칸'' 장치가
쓸 자리일 수 있다.
\item \textbf{웹툰도 $-0.0204$ 다} --- 같은 방향인데 크기가 $1/3$ 이고
시계열은 여전히 값지다($-0.0539$). \textbf{문턱이 어디인가.}
\end{itemize}

\end{document}
